\subsection{Affine harmonic sequences as signed moments}
\label{affine:sec:moments}
The signed-measure formulation also certifies sequences whose harmonic and
denominator steps differ. The Cayley continuation supplies the integer-order
case; its Gamma-density argument extends without change to real outer order
$a\ge1$. It complements the sharper one-sided bounds above.

\begin{theorem}[Affine signed-moment bound]\label{affine:thm:moments}
Let $r,d$ be positive integers, $a\ge1$ real and $b>0$ real. Set
\[
 A_k=\frac{H_{rk}^{(b)}}{(dk+1)^a},\qquad
 m_1=1,\qquad m_a=\frac{(a-1)^{a-1}e^{-(a-1)}}{\Gamma(a)}\quad(a>1).
\]
There is a finite signed measure $\nu$ on $[0,1]$, with no atom at one,
such that $A_k=\int x^k\,d\nu(x)$ and
\[
 \|\nu\|_{\rm TV}\le\frac{2r}{d}m_a b\zeta(b+1).
\]
For $b>1$ one also has $\|\nu\|_{\rm TV}\le2\zeta(b)$.
Consequently the ordinary alternating series converges, and its $N$-term
Euler transform $E_N$ satisfies
\[
 \left|\sum_{k\ge0}(-1)^kA_k-E_N\right|\le C2^{-N},
\]
where $C$ may be either applicable bound. For integer $b$, valid rational
choices are $C=10r/(3d)$ at $b=1$ and $C=10/3$ at $b\ge2$.
\end{theorem}
\begin{proof}
Let $Y$ have Gamma density $y^{a-1}e^{-y}/\Gamma(a)$, and let $\mu$ be
the law of $X=e^{-dY}$. Then $\int x^k\,d\mu=(dk+1)^{-a}$.
In the logarithmic coordinate $s=-\log x$, its density is
\[
 f(s)=\frac{s^{a-1}e^{-s/d}}{d^a\Gamma(a)}\,\mathbf1_{s>0}.
\]
For every real $a\ge1$ its maximum is $m_a/d$. Extended by zero to the
negative axis, the total variation of its distributional derivative is
$2m_a/d$: the rise and fall give this for $a>1$, while at $a=1$ the
jump at zero and the decreasing exponential each contribute $1/d$.
The $L^1$ distance from a bounded-variation density to its translate by
$h\ge0$ is at most $h$ times this derivative variation. Indeed, integrate
$f(s-h)-f(s)=-\int_{s-h}^s f'(v)\,dv$ and use Fubini; smooth convolution
proves the same bound for the jump case. Thus, writing $D_\rho\mu$ for
pushforward by $x\mapsto\rho x$,
\[
 \|\mu-D_{e^{-rt}}\mu\|_{\rm TV}
 \le\min\left(2,\frac{2m_ar}{d}t\right).
\]
The generalized harmonic integral gives the measure
\[
 \nu=\frac1{\Gamma(b)}\int_0^\infty
   \frac{t^{b-1}}{e^t-1}(\mu-D_{e^{-rt}}\mu)\,dt.
\]
This integral converges in total variation: the translation bound supplies
an additional factor $t$ at zero, and exponential decay controls infinity.
Taking moments gives $A_k$, since the moments of the measure difference are
$(dk+1)^{-a}(1-e^{-rkt})$. Its total-variation bounds follow from
\[
 \frac1{\Gamma(b)}\int_0^\infty\frac{t^b}{e^t-1}\,dt
 =b\zeta(b+1),\qquad
 \frac1{\Gamma(b)}\int_0^\infty\frac{t^{b-1}}{e^t-1}\,dt
 =\zeta(b)\quad(b>1).
\]
Each constituent has no atom at one, and variation domination passes that
property to $\nu$. The finite alternating kernel is bounded by two and
converges for $x<1$ to $(1+x)^{-1}$, so dominated convergence proves
ordinary convergence. The exact remainder kernel in the Euler expansion is
$2^{-N}(1-x)^N/(1+x)$, bounded by $2^{-N}$; integrating against $\nu$
proves the stated error. Finally $m_a\le1$, by bounding the maximum of the
unit-scale Gamma density, and $\zeta(2)<5/3$. For the displayed rational
bound at $b=1$ one may also use only the integer $a$ needed by the rational
evaluator: then $m_a\le1$ follows directly from
$e^{a-1}\ge(a-1)^{a-1}/(a-1)!$. For the real-$a$ inequality, write the
Gamma density as the convolution of an exponential density and a
Gamma$(a-1,1)$ probability density when $a>1$; its supremum is at most one.
For integer $b\ge2$, $\zeta(b)\le\zeta(2)$ gives $10/3$.
\end{proof}
The real-order theorem does not make its centers rational at noninteger
indices. The finite rational implementation uses integer $a,b,r,d$ and
binomial-tail Euler weights. No complete monotonicity of $A_k$ is assumed:
the representing measure is signed and has total mass zero, since $A_0=0$.
